\documentclass{article}
\usepackage{amsmath, amssymb, amsthm}
\usepackage{hyperref}
\usepackage{geometry}
\geometry{margin=1in}

\title{Erd\H{o}s Problem \#626}
\date{Last updated: 2025-08-31}
\author{Source: erdosproblems.com}

\begin{document}
\maketitle

\noindent\textbf{Status:} OPEN \quad \textbf{No prize}

\noindent\textbf{Tags:} graph theory, chromatic number, cycles

\medskip
\noindent\textbf{Problem Statement:}

Let $k\geq 4$ and $g_k(n)$ denote the largest $m$ such that there is a graph on $n$ vertices with chromatic number $k$ and girth $>m$ (i.e. contains no cycle of length $\leq m$). Does\[\lim_{n\to \infty}\frac{g_k(n)}{\log n}\]exist?

Conversely, if $h^{(m)}(n)$ is the maximal chromatic number of a graph on $n$ vertices with girth $>m$ then does\[\lim_{n\to \infty}\frac{\log h^{(m)}(n)}{\log n}\]exist, and what is its value?

\bigskip
\noindent\textbf{Remarks:}

It is known that\[\frac{1}{4\log k}\log n\leq g_k(n) \leq \frac{2}{\log(k-2)}\log n+1,\]the lower bound due to Kostochka \cite{Ko88} and the upper bound to Erd\H{o}s \cite{Er59b}. 

Erd\H{o}s \cite{Er59b} proved that\[\lim_{n\to \infty}\frac{\log h^{(m)}(n)}{\log n}\gg \frac{1}{m}\]and, for odd $m$,\[\lim_{n\to \infty}\frac{\log h^{(m)}(n)}{\log n}\leq \frac{2}{m+1},\]and conjectured this is sharp. He had no good guess for the value of the limit for even $m$, other that it should lie in $[\frac{2}{m+2},\frac{2}{m}]$, but could not prove this even for $m=4$.

See also the entry in the graphs problem collection.

\bigskip
\noindent\textbf{References:}

[Er59b] Erd\H{o}s, P., Graph theory and probability. Canadian J. Math. (1959), 34-38.
 
 [Ko88] Kostochka, A. V., Upper bounds on the chromatic number of graphs. Trudy Inst. Mat. (Novosibirsk) (1988), 204-226, 265.

\bigskip
\noindent\small{Source: \url{https://www.erdosproblems.com/626}}
\end{document}
